The gene expression datasets used are Broad and Sanger.
These were extracted from the ``RNA-Seq 2022-06-24'' data collection,
downloaded from the Cell Model Passports hub~\cite{van_der_meer_cell_2019, noauthor_cell_nodate}.
This collection contains RNA transcript counts of thousands of genes, including PD-L1, for hundreds of samples.
The measurements of each sample were done on cancer cell lines.
These are derived from human tissue and are often immortalized.
The advantage is that cancer cell lines behave more consistently in experiments compared to fresh human cancer tissue.
We refer to a sample as a ``datapoint'' or simply ``point''.
Combined, the datasets contain $1291$ points.
These are too few points to use methods likes neural networks.

For each sample, the tissue status and tissue type is also recorded in the dataset.
The samples consist of a wide variety of cancer types found in various types of tissue.
Tissue status among the samples is not homogenous.
Sample counts per dataset and tissue status are given in \autoref{tab:data_counts}.

\begin{table}[h]
    \centering
    \resizebox{\textwidth}{!}{%
\begin{tabular}{l rrrrr r}
\toprule
\textbf{Dataset} & \textbf{Tumor} & \textbf{Metastasis} & \textbf{Papiloma} & \textbf{Normal} & \textbf{Unknown} & \textbf{Total} \\
\midrule
\textbf{Broad}  & 507 & 272 & 1 & 1 & 68  & 849  \\
\textbf{Sanger} & 240 & 113 & 0 & 6 & 83  & 442  \\
\midrule
\textbf{Total}  & 747 & 385 & 1 & 7 & 151 & 1291 \\
\bottomrule
\end{tabular}
}
    \caption{Sample counts by dataset and tissue type.}
    \label{tab:data_counts}
\end{table}

Out of all candidate TF, a total of 18 were selected to be of interest.
This selection was based on literature, in particular two criteria:
(1) the TF can physically bind to the CD274 gene sequence and (2)
there is a recorded functional response in literature.
A recorded functional response means that some significant change in PD-L1 expression is associated
with the activity of the TF\@.

The TF activities are estimated using DecoupleR model~\cite{badia-i-mompel_decoupler_2022, badia-i-mompel_transcription_nodate},
using CollecTRI as prior knowledge~\cite{muller-dott_expanding_2023}.
DecoupleR estimates TF activity in a sample based on the full gene expression profile (thousands of genes).
Simply put, it estimates TF activity based on the expression (level of under- or overregulation)
of the genes it regulates (according to the prior knowledge).
The output is, for each TF, a ``TF activity score'', also referred to as ``TF activity'' or ``TF activity estimate''.
Descriptive statistics for the corresponding TF activities and PD-L1 expression in both Broad and Sanger
can be found in \autoref{tab:descr_broad} and \autoref{tab:descr_sanger}, respectively.



The cross-correlation matrix of TF activity scores is depicted in \autoref{fig:corr_activity}.
The cross-correlation matrix of the TF gene expressions (used to calculate TF activities) is depicted in \autoref{fig:corr_TPM}.
Notably, TF activity scores are highly correlated.
This indicates that the selected TFs tend to be active or inactive simultaneously.
This may be expected, if these TFs are known to regulate similar genes.
Meanwhile, gene expressions are not highly correlated.
This suggests that gene expression is not an accurate indicator of simultaneous TF activity.
Alternatively, TF activity scores could be biased towards showing correlation due to the method of estimation.
Both figures also include correlation with PD-L1 expression in the bottom row.
The PD-L1 expression shows relatively low correlation with TF activity scores and gene expressions, respectively.
This indicates that activity or expression of a single TF does not explain PD-L1 expression substantially.
This motivates us further to consider the effect of groups of TFs.


PD-L1 expression for Broad and Sanger dataset combined was normalized to mean $0$ and standard deviation $1$,
The histogram of PD-L1 expression of Broad samples is depicted in \autoref{fig:pdl1_expression_broad}.
It shows that the majority of observations clusters around the minimum.
Note that the plot starts at the minimum observed expression value and is capped at an expression value of $0.5$,
while the maximum observed PD-L1 expression is $3.55$.

\begin{figure}
    \centering
    \includesvg[width=\textwidth]{gen/data/pdl1_expression_broad.svg}
    \caption{Histogram of PD-L1 expression in Broad dataset. Increasing quantiles are also depicted from left to right using vertical lines.}
    \label{fig:pdl1_expression_broad}
\end{figure}

The goal is to infer general combinations of TF contributing to PD-L1 expression.
Hence, it is important that the rulesets produced generalize to other datasets.
This means rulesets need to be tested on unseen data.
Broad serves as the training set.
Cross-validation will be used on splits of Broad to study how well the results various methods, under various parameters, generalize.
Sanger serves as the holdout dataset and will be used only once, at the very end (\autoref{ch:holdout}).
This is the ultimate test whether the method is able to produce general rulesets.

%
%\section{Transcription Factor activity}\label{sec:data_TF_activity}
%Gene expression is not equivalent to protein activity.
%If gene expression is high, it does not necessarily mean that the resulting protein is active.
%There are intermediate steps that can affect the protein activity.
%For a TF to bind to a gene, however, the TF needs to be active.
%Therefore, using gene expression directly may not be suitable for our analysis.
%To address this issue, an estimate of TF activities is calculated using
%the decoupleR model with CollecTRI as prior knowledge network
%\cite{badia-i-mompel_decoupler_2022, muller-dott_expanding_2023, badia-i-mompel_transcription_nodate}.
%
%DecoupleR estimates TF activity based on the expression of all genes in the dataset via a collection of linear models. \todo{How much should I go into this model?}
%
%
%

%
%\subsection{Tissue status}\label{subsec:tissue_status}
